\section{Deformations of dimer models}\label{sec_def_deform}

In this section, we will introduce the concept of a deformation of consistent dimer models.
This operation is defined for type I zigzag paths on a consistent dimer model, and there are two kinds of deformations,
which we call the zig-deformation (see Definition~\ref{def_deformation_zig}) and the zag-deformation (see Definition~\ref{def_deformation_zag}).
For some special cases, these deformations preserve the consistency condition (see Proposition~\ref{prop_preserve_consistent1}),
but they change the associated PM polygon, whereas we will see that the PM polygon of the deformed dimer model is exactly the combinatorial mutation of a~polygon (see Theorem~\ref{mutation=deformation1}).

\subsection{Definition of deformations of dimer models}\label{subsec_def_deform}

Let $\Gamma$ be a reduced consistent dimer model.
In particular, any slope of a zigzag path on $\Gamma$ is primitive.
Let $\calZ_v(\Gamma)$ be the subset of zigzag paths on $\Gamma$ whose slopes are the primitive vector~$v\in\ZZ^2$,
and $\calZ_v^\rmI(\Gamma)$ be the subset of $\calZ_v(\Gamma)$ consisting of type I zigzag paths.
We first prepare the \emph{deformation data}.

\begin{Definition}[deformation data]\label{def_deformation_data}
Let $\Gamma$ be a reduced consistent dimer model. In order to define the deformation of $\Gamma$, we fix the following data.
\begin{itemize}\itemsep=0pt
\item[(1)] We choose a type I zigzag path $z$, and let $2n\coloneqq\ell(z)$ and $v\coloneqq [z]$.
\item[(2)] We then fix a positive integer $r$ such that $r\le |\calZ_v^\rmI(\Gamma)|$, and let $h=n-r$.
\item[(3)] We take a subset $\{z_1,\dots,z_r\}\subset\calZ_v^\rmI(\Gamma)$ of type I zigzag paths, in which case we have $2n=\ell(z_1)=\cdots=\ell(z_r)$ by Lemma~\ref{lem_same_length}.
Therefore, each $z_i$ can be described as
\begin{displaymath} z_i=z_i[1]z_i[2]\cdots z_i[2n-1]z_i[2n]. \end{displaymath}
\item[(4)]
We consider non-negative integers $p_1,\dots,p_r\in\ZZ_{\ge0}$ such that $p_1+\cdots+p_r=n-r=h$.
We call $\bfp\coloneqq(p_1,\dots,p_r)$ a \emph{deformation weight} of $\{z_1, \dots, z_r\}$.
\end{itemize}
\end{Definition}

\begin{Remark}\label{rem_typeI_to_typeII}
To define the deformation data, we need a type I zigzag path.
If a dimer model is isoradial, then any zigzag path is type I (see Definition~\ref{def_isoradial}), and hence $\big|\calZ_v^\rmI(\Gamma)\big|=|\calZ_v(\Gamma)|$.
Also, even if $\Gamma$ contains no type I zigzag paths, we sometimes change a type II zigzag path into type I by using the \emph{mutations of dimer models}
(see Appendix~\ref{app_mutationdimer}, especially Example~\ref{mutation_typeII_to_I}).
\end{Remark}

\begin{Definition}[zig-deformation]\label{def_deformation_zig}
Let the notation be the same as in Definition~\ref{def_deformation_data}.
For a~deformation weight $\bfp=(p_1,\dots,p_r)$ of $\{z_1, \dots, z_r\}$, we consider the following procedures:
\begin{enumerate}[(\text{zig}-1)]
\setlength{\parskip}{0pt}\itemsep=0pt
\item Using split moves, we insert $p_i$ white nodes and $p_i$ black nodes in each zig of $z_i$.

\noindent{\bf [Notation]}
\begin{itemize}\itemsep=0pt
\item For any zig $z_i[2m-1]$ of $z_i$ where $m=1,\dots,n$ and $i=1,\dots,r$, we denote by $b_i[2m-1]$ (resp.~$w_i[2m-1]$) the black (resp.\ white) node that is the endpoint of $z_i[2m-1]$.
\item We denote the white nodes added in the zig $z_i[2m-1]$ by $w_{i,1}[2m-1],\dots,w_{i,p_i}[2m-1]$, and the black ones by $b_{i,1}[2m-1],\dots,b_{i,p_i}[2m-1]$.
Here, the subscripts increase in the direction from $b_i[2m-1]$ to $w_i[2m-1]$.
\end{itemize}
\item We remove all zags of $z_i$ for every $i=1,\dots,r$.
\item If $p_i\neq 0$, then we connect the white node $w_{i,j}[2m-1]$ to the black node $b_{i,j}[2m+1]$ where $j=1,\dots,p_i$ and $m=1,\dots,n$.
(Note that $w_{i,j}[2n-1]$ is connected to $b_{i,j}[2n+1]\coloneqq b_{i,j}[1]$.)
We denote by $z_{i,j}$ the new $1$-cycle, which will be a zigzag path on the deformed dimer model, obtained by connecting
\begin{displaymath}
w_{i,j}[2n-1], b_{i,j}[2n-1], w_{i,j}[2n-3], b_{i,j}[2n_i-3],\dots,w_{i,j}[1], b_{i,j}[1]
\end{displaymath}
cyclically ($i=1,\dots,r$ and $j=1,\dots,p_i$).
\item[(join)] If there exist $2$-valent nodes, then we apply the join moves to the dimer model obtained by the above procedures and make it reduced.
\end{enumerate}
We denote the resulting dimer model by $\nu^\zig_\bfp(\Gamma, \{z_1,\dots,z_r\})$,
and call it the \emph{zig-deformation of~$\Gamma$ at $\{z_1,\dots,z_r\}$ with respect to the weight $\bfp$}.
If a situation is clear, we simply denote this by~$\nu^\zig_\bfp(\Gamma)$.
\end{Definition}

\begin{figure}[h!]\centering
\scalebox{0.55}{
\begin{tikzpicture}
\newcommand{\edgewidth}{0.05cm} %line_width_of _edges
\newcommand{\nodewidth}{0.05cm} %line_width_around_nodes
\newcommand{\noderad}{0.16} %radius_of_node
%coordinate setting
\coordinate (W1) at (0,1.2); \coordinate (W2) at (0,3.6);
%\coordinate (B1) at (3,0); \coordinate (B2) at (3,2.4);
\path (W1) ++(-20:4cm) coordinate (B1); \path (W2) ++(-20:4cm) coordinate (B2);
\path (B1) ++(200:2cm) coordinate (B0); \path (W2) ++(20:2cm) coordinate (W3);

\path (W1) ++(-20:0.8cm) coordinate (B1-1); \path (W1) ++(-20:1.6cm) coordinate (W1-1);
\path (W1) ++(-20:2.4cm) coordinate (B1-2); \path (W1) ++(-20:3.2cm) coordinate (W1-2);
\path (W2) ++(-20:0.8cm) coordinate (B2-1); \path (W2) ++(-20:1.6cm) coordinate (W2-1);
\path (W2) ++(-20:2.4cm) coordinate (B2-2); \path (W2) ++(-20:3.2cm) coordinate (W2-2);

\path (B1) ++(30:1cm) coordinate (B1e); \path (B1) ++(330:1cm) coordinate (B1s);
\path (W1) ++(150:1cm) coordinate (W1w); \path (W1) ++(210:1cm) coordinate (W1s);
\path (B2) ++(30:1cm) coordinate (B2e); \path (B2) ++(330:1cm) coordinate (B2s);
\path (W2) ++(150:1cm) coordinate (W2w); \path (W2) ++(210:1cm) coordinate (W2s);

\path (W1) ++(160:0.7cm) coordinate (W1+); \path (W1) ++(200:0.7cm) coordinate (W1-);
\path (W2) ++(160:0.7cm) coordinate (W2+); \path (W2) ++(200:0.7cm) coordinate (W2-);
\path (B1) ++(20:0.7cm) coordinate (B1+); \path (B1) ++(340:0.7cm) coordinate (B1-);
\path (B2) ++(20:0.7cm) coordinate (B2+); \path (B2) ++(340:0.7cm) coordinate (B2-);

\node (original) at (0,0) {
\begin{tikzpicture}
%edge
\draw [line width=\edgewidth] (B1)--(W1); \draw [line width=\edgewidth] (B2)--(W1) ; \draw [line width=\edgewidth] (B2)--(W2) ;
\draw [line width=\edgewidth] (B1)--(B0) ; \draw [line width=\edgewidth] (W2)--(W3) ;
\draw [line width=\edgewidth] (B1)--(B1e); \draw [line width=\edgewidth] (B1)--(B1s);
\draw [line width=\edgewidth] (W1)--(W1w); \draw [line width=\edgewidth] (W1)--(W1s);
\draw [line width=\edgewidth] (B2)--(B2e); \draw [line width=\edgewidth] (B2)--(B2s);
\draw [line width=\edgewidth] (W2)--(W2w); \draw [line width=\edgewidth] (W2)--(W2s);

\draw [line width=\edgewidth,line cap=round, dash pattern=on 0pt off 2.5\pgflinewidth] (W1+)--(W1-);
\draw [line width=\edgewidth,line cap=round, dash pattern=on 0pt off 2.5\pgflinewidth] (W2+)--(W2-);
\draw [line width=\edgewidth,line cap=round, dash pattern=on 0pt off 2.5\pgflinewidth] (B1+)--(B1-);
\draw [line width=\edgewidth,line cap=round, dash pattern=on 0pt off 2.5\pgflinewidth] (B2+)--(B2-);
%blacknode
\draw [line width=\nodewidth, fill=black] (B1) circle [radius=\noderad] ; \draw [line width=\nodewidth, fill=black] (B2) circle [radius=\noderad] ;
%whitenode
\draw [line width=\nodewidth, fill=white] (W1) circle [radius=\noderad] ; \draw [line width=\nodewidth, fill=white] (W2) circle [radius=\noderad] ;

%zigzag path
\coordinate (W1z) at (0,1.2); \coordinate (W2z) at (0,3.6);
\path (W1) ++(90:0.2cm) coordinate (W1z); \path (W2) ++(90:0.2cm) coordinate (W2z);
\path (W1z) ++(-20:4cm) coordinate (B1z); \path (W2z) ++(-20:4cm) coordinate (B2z);
\path (B1z) ++(200:2cm) coordinate (B0z); \path (W2z) ++(20:2cm) coordinate (W3z);
\path (B1z) ++(200:2cm) coordinate (B0z); \path (W2z) ++(20:2.2cm) coordinate (W3z);

\path (W1) ++(-90:0.2cm) coordinate (W1zz); \path (W1zz) ++(150:1.5cm) coordinate (W1wzz);
\draw [->, rounded corners, line width=0.08cm, red] (B0z)--(B1z)--(W1z)--(B2z)--(W2z)--(W3z);

\node at (2.5,3.5) {\color{red}\large$z_i[2m+1]$};
\node at (1.2,2.2) {\color{red}\large$z_i[2m]$};
\node at (2.5,1.1) {\color{red}\large$z_i[2m-1]$};
\end{tikzpicture}};

\node (original_insert) at (10,0) {
\begin{tikzpicture}
%edge
\draw [line width=\edgewidth] (B1)--(W1); \draw [line width=\edgewidth] (B2)--(W1) ; \draw [line width=\edgewidth] (B2)--(W2) ;
\draw [line width=\edgewidth] (B1)--(B0) ; \draw [line width=\edgewidth] (W2)--(W3) ;
\draw [line width=\edgewidth] (B1)--(B1e); \draw [line width=\edgewidth] (B1)--(B1s);
\draw [line width=\edgewidth] (W1)--(W1w); \draw [line width=\edgewidth] (W1)--(W1s);
\draw [line width=\edgewidth] (B2)--(B2e); \draw [line width=\edgewidth] (B2)--(B2s);
\draw [line width=\edgewidth] (W2)--(W2w); \draw [line width=\edgewidth] (W2)--(W2s);

\draw [line width=\edgewidth,line cap=round, dash pattern=on 0pt off 2.5\pgflinewidth] (W1+)--(W1-);
\draw [line width=\edgewidth,line cap=round, dash pattern=on 0pt off 2.5\pgflinewidth] (W2+)--(W2-);
\draw [line width=\edgewidth,line cap=round, dash pattern=on 0pt off 2.5\pgflinewidth] (B1+)--(B1-);
\draw [line width=\edgewidth,line cap=round, dash pattern=on 0pt off 2.5\pgflinewidth] (B2+)--(B2-);
%blacknode
\draw [line width=\nodewidth, fill=black] (B1) circle [radius=\noderad] ; \draw [line width=\nodewidth, fill=black] (B2) circle [radius=\noderad] ;
\draw [line width=\nodewidth, fill=black] (B1-1) circle [radius=\noderad] ; \draw [line width=\nodewidth, fill=black] (B1-2) circle [radius=\noderad] ;
\draw [line width=\nodewidth, fill=black] (B2-1) circle [radius=\noderad] ; \draw [line width=\nodewidth, fill=black] (B2-2) circle [radius=\noderad] ;
%whitenode
\draw [line width=\nodewidth, fill=white] (W1) circle [radius=\noderad] ; \draw [line width=\nodewidth, fill=white] (W2) circle [radius=\noderad] ;
\draw [line width=\nodewidth, fill=white] (W1-1) circle [radius=\noderad] ; \draw [line width=\nodewidth, fill=white] (W1-2) circle [radius=\noderad] ;
\draw [line width=\nodewidth, fill=white] (W2-1) circle [radius=\noderad] ; \draw [line width=\nodewidth, fill=white] (W2-2) circle [radius=\noderad] ;
\end{tikzpicture}};


\node (original_remove) at (10,-6.5) {
\begin{tikzpicture}
%edge
\draw [line width=\edgewidth] (B1)--(W1); \draw [line width=\edgewidth] (B2)--(W2);
%\draw [line width=\edgewidth] (B2)--(W1); \draw [line width=\edgewidth] (B1)--(B0); \draw [line width=\edgewidth] (W2)--(W3);
\draw [line width=\edgewidth] (B1)--(B1e); \draw [line width=\edgewidth] (B1)--(B1s);
\draw [line width=\edgewidth] (W1)--(W1w); \draw [line width=\edgewidth] (W1)--(W1s);
\draw [line width=\edgewidth] (B2)--(B2e); \draw [line width=\edgewidth] (B2)--(B2s);
\draw [line width=\edgewidth] (W2)--(W2w); \draw [line width=\edgewidth] (W2)--(W2s);

\draw [line width=\edgewidth,line cap=round, dash pattern=on 0pt off 2.5\pgflinewidth] (W1+)--(W1-);
\draw [line width=\edgewidth,line cap=round, dash pattern=on 0pt off 2.5\pgflinewidth] (W2+)--(W2-);
\draw [line width=\edgewidth,line cap=round, dash pattern=on 0pt off 2.5\pgflinewidth] (B1+)--(B1-);
\draw [line width=\edgewidth,line cap=round, dash pattern=on 0pt off 2.5\pgflinewidth] (B2+)--(B2-);
%blacknode
\draw [line width=\nodewidth, fill=black] (B1) circle [radius=\noderad] ; \draw [line width=\nodewidth, fill=black] (B2) circle [radius=\noderad] ;
\draw [line width=\nodewidth, fill=black] (B1-1) circle [radius=\noderad] ; \draw [line width=\nodewidth, fill=black] (B1-2) circle [radius=\noderad] ;
\draw [line width=\nodewidth, fill=black] (B2-1) circle [radius=\noderad] ; \draw [line width=\nodewidth, fill=black] (B2-2) circle [radius=\noderad] ;
%whitenode
\draw [line width=\nodewidth, fill=white] (W1) circle [radius=\noderad] ; \draw [line width=\nodewidth, fill=white] (W2) circle [radius=\noderad] ;
\draw [line width=\nodewidth, fill=white] (W1-1) circle [radius=\noderad] ; \draw [line width=\nodewidth, fill=white] (W1-2) circle [radius=\noderad] ;
\draw [line width=\nodewidth, fill=white] (W2-1) circle [radius=\noderad] ; \draw [line width=\nodewidth, fill=white] (W2-2) circle [radius=\noderad] ;
\end{tikzpicture}};


\node (before) at (20,-6.5) {
\begin{tikzpicture}
%edge
\draw [line width=\edgewidth] (B1)--(W1); \draw [line width=\edgewidth] (B2)--(W2) ;
\draw [line width=\edgewidth] (B1)--(B1e); \draw [line width=\edgewidth] (B1)--(B1s);
\draw [line width=\edgewidth] (W1)--(W1w); \draw [line width=\edgewidth] (W1)--(W1s);
\draw [line width=\edgewidth] (B2)--(B2e); \draw [line width=\edgewidth] (B2)--(B2s);
\draw [line width=\edgewidth] (W2)--(W2w); \draw [line width=\edgewidth] (W2)--(W2s);

\draw [line width=\edgewidth,line cap=round, dash pattern=on 0pt off 2.5\pgflinewidth] (W1+)--(W1-);
\draw [line width=\edgewidth,line cap=round, dash pattern=on 0pt off 2.5\pgflinewidth] (W2+)--(W2-);
\draw [line width=\edgewidth,line cap=round, dash pattern=on 0pt off 2.5\pgflinewidth] (B1+)--(B1-);
\draw [line width=\edgewidth,line cap=round, dash pattern=on 0pt off 2.5\pgflinewidth] (B2+)--(B2-);

\draw [line width=\edgewidth] (B2-1)--(W1-1); \draw [line width=\edgewidth] (B2-2)--(W1-2);
\path (B1-1) ++(-70:1.5cm) coordinate (B1-1-); \draw [line width=\edgewidth] (B1-1)--(B1-1-);
\path (B1-2) ++(-70:1.5cm) coordinate (B1-2-); \draw [line width=\edgewidth] (B1-2)--(B1-2-);
\path (W2-1) ++(110:1.5cm) coordinate (W2-1-); \draw [line width=\edgewidth] (W2-1)--(W2-1-);
\path (W2-2) ++(110:1.5cm) coordinate (W2-2-); \draw [line width=\edgewidth] (W2-2)--(W2-2-);
%blacknode
\draw [line width=\nodewidth, fill=black] (B1) circle [radius=\noderad] ; \draw [line width=\nodewidth, fill=black] (B2) circle [radius=\noderad] ;
\draw [line width=\nodewidth, fill=black] (B1-1) circle [radius=\noderad] ; \draw [line width=\nodewidth, fill=black] (B1-2) circle [radius=\noderad] ;
\draw [line width=\nodewidth, fill=black] (B2-1) circle [radius=\noderad] ; \draw [line width=\nodewidth, fill=black] (B2-2) circle [radius=\noderad] ;
%whitenode
\draw [line width=\nodewidth, fill=white] (W1) circle [radius=\noderad] ; \draw [line width=\nodewidth, fill=white] (W2) circle [radius=\noderad] ;
\draw [line width=\nodewidth, fill=white] (W1-1) circle [radius=\noderad] ; \draw [line width=\nodewidth, fill=white] (W1-2) circle [radius=\noderad] ;
\draw [line width=\nodewidth, fill=white] (W2-1) circle [radius=\noderad] ; \draw [line width=\nodewidth, fill=white] (W2-2) circle [radius=\noderad] ;
\end{tikzpicture}};

\iffalse
\node (after) at (20,-6.5) {
\begin{tikzpicture}
%edge
\draw [line width=\edgewidth] (B1)--(W1); \draw [line width=\edgewidth] (B2)--(W2) ;
\draw [line width=\edgewidth] (B1)--(B1e); \draw [line width=\edgewidth] (B1)--(B1s);
\draw [line width=\edgewidth] (W1)--(W1w); \draw [line width=\edgewidth] (W1)--(W1s);
\draw [line width=\edgewidth] (B2)--(B2e); \draw [line width=\edgewidth] (B2)--(B2s);
\draw [line width=\edgewidth] (W2)--(W2w); \draw [line width=\edgewidth] (W2)--(W2s);

\draw [line width=\edgewidth,line cap=round, dash pattern=on 0pt off 2.5\pgflinewidth] (W1+)--(W1-);
\draw [line width=\edgewidth,line cap=round, dash pattern=on 0pt off 2.5\pgflinewidth] (W2+)--(W2-);
\draw [line width=\edgewidth,line cap=round, dash pattern=on 0pt off 2.5\pgflinewidth] (B1+)--(B1-);
\draw [line width=\edgewidth,line cap=round, dash pattern=on 0pt off 2.5\pgflinewidth] (B2+)--(B2-);

\draw [line width=\edgewidth] (B2-1)--(W1-1); \draw [line width=\edgewidth] (B2-2)--(W1-2);
\path (B1-1) ++(-70:1.5cm) coordinate (B1-1-); \draw [line width=\edgewidth] (B1-1)--(B1-1-);
\path (B1-2) ++(-70:1.5cm) coordinate (B1-2-); \draw [line width=\edgewidth] (B1-2)--(B1-2-);
\path (W2-1) ++(110:1.5cm) coordinate (W2-1-); \draw [line width=\edgewidth] (W2-1)--(W2-1-);
\path (W2-2) ++(110:1.5cm) coordinate (W2-2-); \draw [line width=\edgewidth] (W2-2)--(W2-2-);

\draw [line width=\edgewidth] (B2-1)--(W1); \draw [line width=\edgewidth] (B2-2)--(W1-1); \draw [line width=\edgewidth] (B2)--(W1-2);
%blacknode
\draw [line width=\nodewidth, fill=black] (B1) circle [radius=\noderad] ; \draw [line width=\nodewidth, fill=black] (B2) circle [radius=\noderad] ;
\draw [line width=\nodewidth, fill=black] (B1-1) circle [radius=\noderad] ; \draw [line width=\nodewidth, fill=black] (B1-2) circle [radius=\noderad] ;
\draw [line width=\nodewidth, fill=black] (B2-1) circle [radius=\noderad] ; \draw [line width=\nodewidth, fill=black] (B2-2) circle [radius=\noderad] ;
%whitenode
\draw [line width=\nodewidth, fill=white] (W1) circle [radius=\noderad] ; \draw [line width=\nodewidth, fill=white] (W2) circle [radius=\noderad] ;
\draw [line width=\nodewidth, fill=white] (W1-1) circle [radius=\noderad] ; \draw [line width=\nodewidth, fill=white] (W1-2) circle [radius=\noderad] ;
\draw [line width=\nodewidth, fill=white] (W2-1) circle [radius=\noderad] ; \draw [line width=\nodewidth, fill=white] (W2-2) circle [radius=\noderad] ;
\end{tikzpicture}};
\fi

\path (original) ++(0:3.5cm) coordinate (original+); \path (original_insert) ++(180:3.5cm) coordinate (original_insert+);
\draw [->, line width=\edgewidth] (original+)--(original_insert+) node[midway,xshift=0cm,yshift=0.5cm] {\LARGE(zig-1)} ;

\path (original_insert) ++(0:3.5cm) coordinate (original_insert-); \path (original_remove) ++(180:3.5cm) coordinate (original_remove+);
\path (original_remove) ++(0:3.5cm) coordinate (original_remove++);
\path (before) ++(180:3.5cm) coordinate (before-); \path (0,-6.5) ++(0:3.5cm) coordinate (before--);

\draw [->, line width=\edgewidth] (before--)--(original_remove+) node[midway,xshift=0cm,yshift=0.5cm] {\LARGE(zig-2)} ;
\draw [->, line width=\edgewidth] (original_remove++)--(before-) node[midway,xshift=0cm,yshift=0.5cm] {\LARGE(zig-3)} ;

%\path (before) ++(0:3.5cm) coordinate (before+); \path (after) ++(180:3.5cm) coordinate (after+);
%\draw [->, line width=\edgewidth] (before+)--(after+) node[midway,xshift=0cm,yshift=0.5cm] {\LARGE(zig-4)} ;
\end{tikzpicture}
}
\caption{The zig-deformation at $z_i$ with $p_i=2$.}\label{fig_zigdeform_p2}
\end{figure}
